\documentclass[12pt]{article}

\usepackage[T1]{fontenc}
\usepackage[utf8]{inputenc}
\usepackage[spanish]{babel}
\usepackage{amsmath,amssymb}
\usepackage{cancel}

\title{Condiciones cinemáticas RHEED}
\author{}
\date{}

\begin{document}
\maketitle

La esfera tiene radio \(\lvert\vec{K}\rvert\), ahora los rods viven en un
plano que corta la esfera, el detalle es saber que en general la ecuación
a resolver es:
\begin{equation}
  \vec{K}-\vec{K}'=\vec{g}_{h,k}+\vec{\lambda}_{k}.
\end{equation}

Las condiciones son que
\begin{equation}
  \lvert\vec{K}\rvert=\lvert\vec{K}'\rvert,
\end{equation}
\begin{equation}
  \lvert\vec{K}\rvert=\lvert\vec{K}-\vec{g}_{h,k}+\vec{\lambda}_{k}\rvert.
\end{equation}

Las condiciones son que
\begin{equation}
  \lvert\vec{k}\rvert=\lvert\vec{k}'\rvert,
\end{equation}
\begin{equation}
  \lvert\vec{k}\rvert=\lvert\vec{k}-\vec{g}_{h,k}+\vec{\lambda}_{k}\rvert,
\end{equation}
\begin{equation}
  \bigl[(\vec{k}\cdot\vec{k})\bigr]^{1/2}
  =
  \bigl[
    (\vec{k}-\vec{g}_{h,k}+\vec{\lambda}_{k})
    \cdot
    (\vec{k}-\vec{g}_{h,k}+\vec{\lambda}_{k})
  \bigr]^{1/2}.
\end{equation}

\begin{equation}
  \lvert\vec{k}\rvert^{2}
  =
  \lvert\vec{k}\rvert^{2}
  -2\lvert\vec{g}_{h,k}\rvert\,\lvert\vec{k}\rvert\cos\theta
  +\lvert\vec{g}_{h,k}\rvert^{2}
  +\lambda_{k}^{2}.
\end{equation}

Cancelando \(\lvert\vec{k}\rvert^{2}\) en ambos lados:
\begin{equation}
  \cancel{\lvert\vec{k}\rvert^{2}}
  =
  \cancel{\lvert\vec{k}\rvert^{2}}
  -2\lvert\vec{g}_{h,k}\rvert\,\lvert\vec{k}\rvert\cos\theta
  +\lvert\vec{g}_{h,k}\rvert^{2}
  +\lambda_{k}^{2},
\end{equation}
\begin{equation}
  \lvert\lambda_{k}\rvert^{2}
  -2\lvert\vec{g}_{h,k}\rvert\,\lvert\vec{k}\rvert\cos\theta
  +\lvert\vec{g}_{h,k}\rvert^{2}
  =0.
\end{equation}

\end{document}
